% ============================================================
%  Literature Review — Direction 1  (10-paper version)
%  Multi-View Mammographic Breast Cancer Detection with
%  Bilateral Asymmetry Fusion and Conformal Uncertainty
%  Quantification
%
%  ACM 40960 · Project 8 — Disease Modelling
%  University College Dublin
%
%  Citation style : APA 7th Edition (biblatex-apa)
%  Font           : Times New Roman (fontspec / XeLaTeX)
%  Spacing        : Double (setspace)
%  Page size      : A4, 1-inch margins
%
%  Overleaf settings:
%    Menu -> Compiler -> XeLaTeX
%    (Biber runs automatically with biblatex under XeLaTeX)
% ============================================================

\documentclass[12pt, a4paper]{article}

% -- Encoding & fonts (XeLaTeX) ------------------------------
\usepackage{fontspec}
\setmainfont{Times New Roman}
\usepackage{unicode-math}
\setmathfont{XITS Math}

% -- Page geometry -------------------------------------------
\usepackage[
  a4paper,
  top    = 2.54cm,
  bottom = 2.54cm,
  left   = 2.54cm,
  right  = 2.54cm
]{geometry}

% -- Line spacing --------------------------------------------
\usepackage{setspace}
\doublespacing

% -- Running header / footer ---------------------------------
\usepackage{fancyhdr}
\setlength{\headheight}{14.5pt}
\pagestyle{fancy}
\fancyhf{}
\fancyhead[L]{\small\textit{Literature Review --- Mammographic Breast Cancer Detection}}
\fancyhead[R]{\small\thepage}
\renewcommand{\headrulewidth}{0.4pt}

% -- Section heading style (APA: bold, flush left) -----------
\usepackage{titlesec}
\titleformat{\section}{\normalfont\bfseries\large}{}{0em}{}
\titleformat{\subsection}{\normalfont\bfseries}{}{0em}{}
\titleformat{\subsubsection}{\normalfont\bfseries\itshape}{}{0em}{}
\titlespacing*{\section}{0pt}{24pt}{6pt}
\titlespacing*{\subsection}{0pt}{18pt}{4pt}

% -- Bibliography (APA 7th via biblatex-apa) -----------------
\usepackage[
  backend   = biber,
  style     = apa,
  sortcites = true,
  sorting   = nyt
]{biblatex}
\addbibresource{litreview_refs_v2.bib}

% -- Hyperlinks ----------------------------------------------
\usepackage[
  colorlinks = true,
  linkcolor  = black,
  citecolor  = black,
  urlcolor   = blue
]{hyperref}
\usepackage{url}

% -- Misc ----------------------------------------------------
\usepackage{microtype}
\setlength{\parindent}{0.5in}
\setlength{\parskip}{0pt}

% -- Abstract environment ------------------------------------
\renewenvironment{abstract}{%
  \begin{center}\textbf{Abstract}\end{center}%
  \begin{quotation}%
}{\end{quotation}}

% ============================================================
\begin{document}

% -- Title page ----------------------------------------------
\begin{titlepage}
  \centering
  \vspace*{3cm}

  {\LARGE\bfseries Literature Review\par}
  \vspace{1.2cm}

  {\large\bfseries
    Multi-View Mammographic Breast Cancer Detection with\\
    Bilateral Asymmetry Fusion and Conformal\\
    Uncertainty Quantification\par}
  \vspace{1.8cm}

  {\normalsize\itshape ACM 40960 \textperiodcentered\ Project 8 --- Disease Modelling\par}
  \vspace{0.6cm}

  {\normalsize University College Dublin\\
  School of Mathematics and Statistics\par}
  \vspace{0.6cm}

  {\normalsize Supervised by Dr.\ Sarp Ak{\c{c}}ay\par}
  \vspace{0.6cm}

  {\normalsize June 2026\par}
\end{titlepage}

% -- Abstract ------------------------------------------------
\begin{abstract}
This literature review surveys the research foundations underpinning a novel
deep learning system for breast cancer detection from digital mammography.
The proposed system makes two primary methodological contributions: multi-view
bilateral asymmetry fusion, implemented via a Siamese encoder with gated
cross-view attention and an explicit asymmetry module
$\delta = |f_L - f_R|$; and conformal uncertainty quantification, providing
distribution-free finite-sample coverage guarantees via Regularised Adaptive
Prediction Sets (RAPS).
The review covers ten primary sources organised across four thematic areas:
the backbone architecture and fine-tuning paradigm that underpin the feature
extractor; the mammography AI literature that establishes clinical credibility
and defines the experimental baselines; the multi-view and bilateral asymmetry
methods that constitute the immediate prior work for the first contribution;
and the conformal prediction literature that grounds the second.
Each paper is assessed for its direct methodological relevance to the proposed
system, with explicit attention to what the proposed work extends, what it
must differentiate itself from, and what it implements directly.
\end{abstract}

\newpage
\tableofcontents
\newpage

% ============================================================
\section{Introduction}

Breast cancer is the most commonly diagnosed cancer among women worldwide.
Mammographic screening programmes reduce mortality through early detection, yet
sensitivity is constrained by reader variability, radiologist fatigue, and the
difficulty of identifying subtle lesions in radiographically dense tissue.
Deep convolutional neural networks have demonstrated radiologist-level
diagnostic performance on large mammography datasets
\parencite{McKinney2020, Wu2019}, catalysing substantial clinical and research
interest in automated screening.

The proposed system targets two underaddressed limitations in the current
literature. First, most high-performing systems process views independently or
fuse them implicitly, discarding the well-established clinical value of
bilateral asymmetry --- the systematic structural difference between left and
right breasts that is one of the most reliable radiological indicators of
malignancy. Second, predictive uncertainty in deployed mammography AI is
routinely characterised by heuristic methods lacking formal guarantees, a
critical shortcoming in a clinical decision-support context where the cost
of overconfident errors is asymmetric and severe.

The proposed system addresses these limitations through a Siamese bilateral
encoder with gated cross-view attention, an explicit asymmetry feature
$\delta = |f_L - f_R|$, and split conformal prediction via RAPS
\parencite{Romano2020}, providing marginal coverage guarantees at any
user-specified miscoverage level $\alpha$.
This review surveys the ten primary sources that establish the intellectual
context, methodological precedents, and comparative baselines for the system.
The review is structured in four thematic sections, each corresponding to a
load-bearing component of the proposed architecture.
All citations follow APA 7th Edition, as standard for quantitative and
computational research submissions at University College Dublin.

% ============================================================
\section{Backbone Architecture and the Fine-Tuning Paradigm}

\subsection{DenseNet: The Architectural Foundation}

The proposed system employs DenseNet-121 as its primary feature extraction
backbone.
\textcite{Huang2017} introduced dense connectivity, in which each layer
receives the concatenated feature maps of \textit{all} preceding layers as
input, rather than only the immediately preceding layer as in residual
networks.
This connectivity pattern promotes feature reuse, substantially reduces
parameter count relative to comparably deep networks, and produces implicit
deep supervision throughout the network by shortening the gradient path from
any layer to the loss.
DenseNet-121 --- the 121-layer configuration --- achieves strong performance
on ImageNet at a parameter budget well-suited to fine-tuning on datasets of
limited size, making it the natural choice for CBIS-DDSM
\parencite{Lee2017}, which contains approximately 2,620 cases.
Understanding the dense connectivity structure is a prerequisite for
reasoning about feature map dimensions and skip connections in the bilateral
fusion module, where the concatenation of encoder outputs from both branches
is combined with the asymmetry feature $\delta$.

\subsection{CheXNet and the Medical Fine-Tuning Template}

\textcite{Rajpurkar2017} fine-tuned DenseNet-121 on the ChestX-ray14 dataset
of 112,120 frontal chest radiographs, achieving F1 scores exceeding the
mean radiologist performance on pneumonia detection.
Their protocol --- ImageNet-pretrained initialisation, aggressive data
augmentation (random horizontal flipping, random crop), and Adam
optimisation with a cyclically decayed learning rate --- established the
canonical transfer learning template for medical imaging with DenseNet.
This protocol is adopted directly in the proposed system's training pipeline.
CheXNet also introduced the practice of reporting performance at the
per-finding level with radiologist comparison, a reporting standard that
the proposed system's evaluation follows for its CBIS-DDSM results.

% ============================================================
\section{Mammography AI: Clinical Baselines}

\subsection{DeepMind International Evaluation}

\textcite{McKinney2020} evaluated a proprietary Google DeepMind mammography
system on 29,426 retrospective cases from the United Kingdom and United
States, supplemented by a prospective UK cohort.
The system reduced false positive rates by 5.7\% (UK) and 1.2\% (US) and
false negative rates by 9.4\% (UK) and 2.7\% (US) relative to a single
radiologist, outperforming all six radiologists who participated in an
independent reader study.
This paper establishes the performance expectations and clinical framing
for the field.
It is the reference every examining supervisor will situate the proposed
work against, and its results define the upper bound of clinical credibility
against which new systems are benchmarked.
The proposed system does not attempt to replicate the scale of this study
--- it trains and evaluates on publicly available data (CBIS-DDSM) --- but
positions its contributions relative to the architectural and methodological
directions this work opened.

\subsection{NYU Deep Learning for Breast Cancer Screening}

\textcite{Wu2019} presented the most methodologically detailed academic
mammography study to date, training on 229,426 screening examinations from
NYU Langone Health.
Their model achieved AUROC 0.876 on an independent test set and produced
statistically significant improvements in radiologist sensitivity when used
as a second reader.
Critically for the proposed system, \textcite{Wu2019} process all four
standard views --- left and right craniocaudal (CC) and mediolateral
oblique (MLO) --- and combine view-level predictions through learned
attention weights.
This is the direct architectural template for the proposed system.
The NYU open-source codebase provides the reference implementation for the
single-view DenseNet-121 baseline against which the bilateral asymmetry
fusion component is evaluated.
The proposed system extends Wu et al.\ by replacing implicit view
combination with a Siamese bilateral encoder, a gated cross-view attention
module, and an explicit asymmetry feature, and by adding conformal
prediction sets at inference time.

\subsection{CBIS-DDSM: The Primary Dataset}

\textcite{Lee2017} describe the Curated Breast Imaging Subset of the Digital
Database for Screening Mammography (CBIS-DDSM), the primary dataset on which
the proposed system is trained and evaluated.
CBIS-DDSM contains 2,620 digitised film-screen mammography cases with expert
annotations for mass and calcification findings, BI-RADS assessment
categories, pathology labels (benign/malignant), bounding boxes, and
pixel-level segmentation masks for a subset of cases.
Reading this paper is a prerequisite before writing any dataloader code.
The dataset's structure has several non-obvious properties that affect
implementation: cases are stored as DICOM files in a non-standard directory
layout; image-level and case-level CSV files must be joined on
\texttt{patient\_id} and laterality to reconstruct bilateral pairs; and the
train/test split is provided at the case level rather than the image level,
requiring careful handling to avoid data leakage between left-right pairs.
The proposed system's bilateral fusion module requires precisely this
patient-level grouping, making the data structure described by
\textcite{Lee2017} architecturally relevant, not merely bibliographically
obligatory.

% ============================================================
\section{Multi-View and Bilateral Asymmetry Fusion}

\subsection{Dual-View Deep Learning for Mammography}

\textcite{Wu2020} extended the single-view NYU architecture into an explicit
dual-view framework, processing CC and MLO views of the same breast jointly
through a shared encoder with a cross-view attention mechanism.
Their system achieved statistically significant improvements over the
single-view baseline on the NYU proprietary dataset, demonstrating that
explicit architectural modelling of the CC/MLO relationship carries
diagnostic information beyond what is captured by independent view processing
with late fusion.
This paper is the most direct architectural prior work for the proposed system
and must be read with particular care.
The key dimension along which the proposed system advances beyond
\textcite{Wu2020} is the \textit{bilateral} axis: their model fuses
views of the same breast (CC and MLO), whereas the proposed system
additionally fuses across the left and right breast, treating the
contralateral pair as a source of discriminative asymmetry signal.
This distinction is the primary claim of novelty and must be stated
precisely in every comparison with this paper.

\subsection{Multi-View Bilateral Asymmetry Detection}

\textcite{Nguyen2019} proposed the most direct computational precursor to the
proposed system's bilateral asymmetry module.
They trained a shared-weight CNN encoder on bilateral mammogram pairs and
computed an explicit asymmetry feature as the absolute difference between
left and right breast representations, $\delta = |f_L - f_R|$, which was
then passed to a classifier jointly with the individual view features.
Their architecture demonstrated that this asymmetry signal improved
classification AUROC by a statistically significant margin over a unilateral
baseline on CBIS-DDSM.
The proposed system directly adopts the $\delta = |f_L - f_R|$
formulation and the shared-weight Siamese encoder structure, extending it
with a gated cross-view attention mechanism that modulates the contribution
of each view representation based on the asymmetry magnitude, and with
the conformal prediction wrapper at inference.

\subsection{Bilateral Asymmetry Counterfactual Network}

\textcite{Yang2021} proposed a bilateral asymmetry guided counterfactual
generating network, in which the asymmetry signal between left and right
breast encodings is used to synthesise a plausible normal counterpart for
an abnormal breast.
Classification is then performed contrastively between the real and
counterfactual representations.
Their system achieves strong performance and demonstrates the diagnostic power
of the bilateral asymmetry signal when exploited as a generative prior.
This paper is included not as a precursor to extend but as a comparison point
to differentiate from.
The proposed system and \textcite{Yang2021} share the core insight that
bilateral asymmetry is a powerful discriminative signal, but differ
fundamentally in mechanism: the proposed system incorporates asymmetry
as a direct discriminative feature fused through attention, preserving
interpretability via standard gradient-based visualisation, whereas
\textcite{Yang2021} use asymmetry as a generative prior, introducing
substantial additional architectural complexity and a generative training
objective.
Being precise about this distinction is necessary to establish the
novelty of the proposed contribution.

% ============================================================
\section{Conformal Prediction and Uncertainty Quantification}

\subsection{Distribution-Free Uncertainty Quantification}

\textcite{Angelopoulos2021} provide the definitive tutorial introduction to
conformal prediction and distribution-free uncertainty quantification,
covering the theoretical foundations, split and full conformal variants,
and the construction of prediction sets with finite-sample coverage
guarantees.
This paper grounds the proposed system's second primary contribution.
Split conformal prediction --- the variant implemented --- partitions
available data into training, calibration, and test sets.
Nonconformity scores are computed on the calibration set, and the
$(1{-}\alpha)(1{+}1/n)$-th quantile of these scores serves as the
threshold for prediction set construction at test time.
The resulting sets satisfy a marginal coverage guarantee: for any
user-specified miscoverage level $\alpha \in (0,1)$, the prediction
set contains the true class with probability at least $1{-}\alpha$,
regardless of the underlying data distribution.
This distribution-free property is the defining advantage of conformal
prediction over calibration approaches such as temperature scaling,
which improve reliability on average but provide no per-instance
statistical guarantee.
In a clinical screening context --- where the consequences of
overconfident false negatives are severe and asymmetric ---
the formal guarantee is not merely a theoretical nicety but a
prerequisite for principled deployment.

\subsection{Regularised Adaptive Prediction Sets (RAPS)}

\textcite{Romano2020} introduced Regularised Adaptive Prediction Sets (RAPS),
extending the basic split conformal framework with a regularisation term that
penalises prediction sets containing many low-probability classes.
Standard split conformal prediction constructs sets of fixed threshold,
which can be arbitrarily large when the model is uncertain.
RAPS addresses this by adding a size-regularisation penalty
$\lambda \cdot \max(0, |\mathcal{C}| - k_{\mathrm{reg}})$ to the
nonconformity score, where $|\mathcal{C}|$ is the current set size and
$k_{\mathrm{reg}}$ is the target size.
The penalty encourages tight prediction sets while preserving the
marginal coverage guarantee.
RAPS is the specific conformal variant implemented in the proposed system,
and its hyperparameters ($\lambda$, $k_{\mathrm{reg}}$) are tuned on the
CBIS-DDSM calibration split.
In the clinical context, the practical consequence is that the system
returns tight prediction sets for confident predictions (typically a
singleton \{benign\} or \{malignant\}) while appropriately expanding
to a two-element set \{benign, malignant\} for ambiguous cases, directly
encoding the clinical recommendation to recall for further assessment.

% ============================================================
\section{Synthesis}

The ten papers reviewed above collectively define the proposed system's
intellectual heritage across four dimensions.

The backbone and fine-tuning papers \parencite{Huang2017, Rajpurkar2017}
establish that DenseNet-121 with ImageNet pre-training is the appropriate
feature extractor for a dataset of CBIS-DDSM's size, and that the
CheXNet training protocol is the correct initialisation strategy.

The clinical baseline papers \parencite{McKinney2020, Wu2019, Lee2017}
define the field's performance expectations, provide the direct
architectural template the proposed system extends, and specify the
dataset whose structure determines the bilateral pairing logic.

The bilateral asymmetry papers \parencite{Wu2020, Nguyen2019, Yang2021}
establish the immediate prior work.
\textcite{Nguyen2019} is the system directly extended;
\textcite{Wu2020} is the closest published architecture;
\textcite{Yang2021} is the paper from which the proposed system
must differentiate itself most carefully.

The conformal prediction papers \parencite{Angelopoulos2021, Romano2020}
provide the theoretical foundation and specific implementation of the
second primary contribution.
Together they establish both why conformal prediction is the appropriate
uncertainty quantification framework in a clinical context (formal
coverage guarantees, distribution-free validity) and exactly which
variant the proposed system implements and why (RAPS for efficient,
size-regularised prediction sets).

No existing system combines bilateral asymmetry fusion, explicit
$\delta$-feature modulation via gated cross-view attention, and
distribution-free conformal uncertainty quantification within a unified
architecture trained on publicly available mammography data.
This combination constitutes the novel contribution the proposed work
makes to the field.

% ============================================================
\section{Conclusion}

This review has surveyed ten primary sources across four thematic areas
directly relevant to the proposed system.
The backbone and transfer learning literature establishes the feature
extractor and training protocol.
The mammography AI literature defines clinical baselines, architectural
templates, and the primary dataset.
The bilateral asymmetry literature provides the immediate precursors that
the proposed system extends and must differentiate from.
The conformal prediction literature grounds the uncertainty quantification
contribution in theory and specifies the implemented method.

The scope of this review is intentionally narrow, covering only those
works that are load-bearing for the system's two primary contributions.
Each paper is cited because it directly informs an architectural decision,
provides a baseline the system must surpass, or must be engaged with
explicitly to establish novelty.
This focus reflects the project's position as a research contribution
grounded in two clearly defined methodological innovations, each with
a well-bounded prior literature.

% ============================================================
\newpage
\printbibliography[title={References}]

\end{document}
